\documentclass[a4paper]{article}
\usepackage{subfiles,geometry,fancyhdr,tikz,longtable,enumitem,setspace}

\newcommand{\ano}{2024} %########## CONFIGURACIÓN DE AÑO

\geometry{
 a4paper,
 landscape,
 left=10mm,
 top=45mm,
 right=10mm,
 bottom=14mm
 }
\renewcommand{\familydefault}
{\sfdefault}
\renewcommand{\headrulewidth}{0pt}
\newcommand{\estrategias}{\setstretch{1.5} \centering \begin{itemize}[leftmargin=0.3cm]\item Estudio dirigido de la bibliografía de archivos digitales seleccionados.\item Aplicación de diferentes métodos de organización de la información tales como mapas y redes conceptuales. \item Desarrollo de trabajos prácticos aplicando contenidos conceptuales a distintas situaciones problemáticas.\item Elaboración de informes de las distintas actividades realizadas.\item Mediación con el profesor por medio de classroom  y digitalización de trabajos\end{itemize}
}
\newcommand{\evaluacion}{\textbf{Formativa} \setstretch{2} Consiste en detector los progresos y las dificultades para hacer los ajustes necesarios en los procesos de enseñanza/aprendizaje   Evaluación continua del seguimiento de las actividades de los alumnos, en los trabajos prácticos y mediación a través de los canales de comunicación adoptado teniendo en cuenta la situación particular del contexto}



\pagestyle{fancy}

\fancyhf{} 

\fancyhead[C]{
    \begin{tikzpicture}[remember picture, overlay]
        \draw  (0,-1+3)  node {\includegraphics{planificaciones/images/image1.png}};
        \draw (0,0.4) node {
\includegraphics{planificaciones/images/image2.png}};
    \end{tikzpicture}
}

\begin{document}

\section*{\centering PLANIFICACIÓN \ano}

\begin{tabular}{r l}
     
\textbf{ÁREA:}&TÉCNICA ESPECÍFICA \\

\textbf{ESPACIO CURRICULAR:} &TERMODINÁMICA Y MÁQUINAS TÉRMICAS. \\

\textbf{CURSO:}&5° 2°\\

\textbf{PROFESOR/A:}& FERREYRA, GUSTAVO DAVID\\

\textbf{ORIENTACIÓN EN:} & ELECTROMECÁNICA  \\
     
\end{tabular}

\subsection*{ACUERDOS INSTITUCIONALES:}
\begin{itemize}
    
    \item Las planificaciones, diagnóstico y programa anual se deben entregar en formato digital al Coordinador de su Área para visar el aspecto disciplinar y sugerir posibles correcciones. Luego éste (COORDINADOR) subirá al Drive del SOE, en la carpeta que corresponda a cada área. 

    \item Cada docente deberá entregar en formato papel al SOE (Asesoría) hasta el 17/4 

    \item Cada área realizará acuerdos para establecer la forma, los tiempos de trabajo y proyectos que se ejecutaran durante el año. 

    \item Se establece 2 reuniones de área anuales obligatorias, 1 en cada periodo. 

    \item Colaborar en el cumplimiento de las Normas de Convivencia Institucional. 

    \item Trabajar las diversas técnicas de estudio durante el primer y segundo año, teniendo en cuenta las profundas dificultades que presentan los alumnos para reconocer ideas principales, realizar resúmenes y gráficos en todos los espacios curriculares, sobre todo en aquellos que requieren más trabajo de lectura, comprensión y producción escrita.  

    \item Seguimiento de logros y dificultades de los alumnos produciendo los reajustes necesarios y aplicando nuevas estrategias metodológicas que permitan modificar el rendimiento de los mismos. 

    \item Generar y unificar criterios, formas de evaluar y canales de comunicación entre los docentes de cada área. 

    \item Trabajo conjunto intra e inter área mediante la formulación de proyectos. 

    \item Aplicar diferentes formatos de evaluación. 

    \item Utilizar estrategias de enseñanza sugerida y variada que se adapten a los alumnos que presentan CUD. 

    \item La carpeta debe contener el programa y el acuerdo áulico. 

    \item Tener presente la heterogeneidad y diversidad del grupo clase. 

    \item Fortalecer las habilidades sociales y los buenos modales, favoreciendo a una convivencia armónica y a un clima de respeto. 

    \item Seleccionar, articular y secuenciar contenidos del área en torno a problemas emergentes, socialmente significativos, vinculados a la salud, al ambiente, debates éticos asociados a temáticas científicas y apelando al plateo de problemas y abordaje de proyectos escolares que partan de situaciones reales. 

    \item Vincular la enseñanza y el aprendizaje con procesos de indagación científica, a partir de la experiencia y/o ideas con la guía docente. 

    \item Integrar las TIC con las propuestas pedagógicas en ciencias. 

    \item Trabajar en proceso de evaluación, integrando el aprendizaje y la enseñanza como un momento continuo de valoración. Considerar el error como un espacio y oportunidad de aprendizaje. 

    \item Realizar las planificaciones en conjunto entre los profesores del mismo espacio curricular acordando saberes. 

    \item Promover la articulación entre los espacios del área Técnica Específica con el resto de las áreas para lograr aprendizajes significativos.  

    \item Preparar a los alumnos para la participación en diferentes muestras y concursos organizados por las instituciones escolares. 


\end{itemize}

\subsection*{ACUERDOS DEL ÁREA}
\begin{itemize}
    \item Ponderación de la carpeta de trabajos prácticos como una nota de proceso más. Se debe consensuar, con los alumnos, el formato y la presentación de la CTP.
    \item Se usará el ABP como recurso para el desarrollo de la planificación. El trabajo conjunto entre distintos espacios curriculares es una opción siempre y cuando sea factible para el docente. 
    \item Se debe entregar a los alumnos un programa con los contenidos de la planificación. Se sugiere que ese documento esté incluido en el contrato aúlico.
    \item En el contrato aúlico se deberán establecer, de forma clara, los criterios de evaluación que se aplicarán durante el desarrollo de la materia. 
    \item Se sugiere la implementación de un calendario que permita organizar y proyectar las actividades de la institución y de la materia en sí. Un ejemplo claro es la determinación de las fechas de la Expo de fin de año. 
\end{itemize}

  
\clearpage
\subfile{tabla.tex}
\subfile{programa.tex}

\end{document}